% appendices.tex — ParBench SC26 Paper Appendices A, B, C
% This file is \input{} from paper.tex — NO preamble, NO \begin{document}
% Figures A.x and B.x are PNGs resolved via \graphicspath to analysis/visualizations/
% Figures C.x are PDFs in figures/ (relative to this directory)

\appendices

%% ============================================================================
%% APPENDIX A: API Selection Rationale
%% ============================================================================

\section{API Selection Rationale}
\label{sec:appendix-a}

This appendix provides the complete quantitative evidence underlying ParBench's
selection of CUDA, OpenMP, and OpenCL as target parallel programming APIs.
Section~\ref{sec:suite-selection} in the main paper presents a condensed 10-API
co-occurrence analysis from the initial 21-repository survey; here we present the
expanded data from the full survey and the kernel-level analysis that informed
the final selection.

\subsection{Full API Co-Occurrence Data}
\label{sec:appendix-a1}

The benchmark landscape survey progressed through multiple iterations, expanding
from an initial seed of 21 repositories to a comprehensive survey of over 80
benchmark suites, mini-applications, proxy applications, libraries, and full
applications. This broader survey captured 30+ distinct parallel programming
APIs, providing a panoramic view of the HPC benchmarking ecosystem.

\begin{figure}[htbp]
  \centering
  \includegraphics[width=\columnwidth]{api_cooccurrence_heatmap_v5.png}
  \caption{Expanded API co-occurrence heatmap across all surveyed repositories.
    Cell values indicate the number of repositories providing implementations in
    both the row and column APIs. The concentrated cluster of high co-occurrence
    among MPI, OpenMP, CUDA, HIP, and OpenACC in the upper-left quadrant reveals
    the core API ecosystem that dominates multi-API benchmark suites. Peripheral
    APIs such as Chapel, STM/TM, HLS, and Python show near-zero co-occurrence,
    reflecting their niche status in HPC benchmarking.}
  \label{fig:api-heatmap}
\end{figure}

Section~\ref{sec:suite-selection} in the main paper discusses the 10-API subset
from the initial 21-repository survey; Figure~\ref{fig:api-heatmap} extends this
to the full expanded survey. The dominant pattern is preserved across survey
scales: a dense co-occurrence cluster centered on MPI, OpenMP, CUDA, and HIP,
with sparser connectivity among directive-based (OpenACC, OpenMP Target),
portability (Kokkos, RAJA, SYCL), and niche APIs.

The force-directed network visualization in Figure~\ref{fig:api-network} makes
the clustering structure explicit. Node sizes are proportional to the number of
benchmarks implementing each API, and edges are drawn only for co-occurrence
counts of 3 or more repositories. The result reveals three natural groupings:

% TODO: Restore or update api_cooccurrence_network_color_v7.png
% (deleted from analysis/visualizations/; a non-v7 version exists in
%  analysis/visualizations/take in paper/api_cooccurrence_network_color.png)
\begin{figure}[htbp]
  \centering
  % \includegraphics[width=\columnwidth]{api_cooccurrence_network_color_v7.png}
  \fbox{\parbox{0.9\columnwidth}{\centering\vspace{1em}\textbf{[Missing: api\_cooccurrence\_network\_color\_v7.png]}\vspace{1em}}}
  \caption{API co-occurrence network. Node size proportional to the number of
    benchmarks implementing each API; edges drawn for co-occurrence $\geq 3$.
    APIs colored by family: vendor GPU (cyan), open standard GPU (red),
    directives (orange), message passing (green), portability libraries (grey),
    CPU parallel libraries (blue), single-source C++ (yellow), PGAS/one-sided
    (pink). The dense central cluster of MPI, OpenMP, CUDA, HIP, OpenACC, and
    OpenCL is clearly separated from peripheral APIs. Isolated nodes (Chapel,
    STM/TM, GSParLib, HLS, Python, Serial) have co-occurrence below~3 with all
    other APIs.}
  \label{fig:api-network}
\end{figure}

\begin{enumerate}
  \item \textbf{Core cluster}: MPI (largest node), OpenMP, CUDA, HIP, OpenACC,
    and OpenCL form a tightly interconnected group. These APIs co-occur
    frequently because the benchmark suites that provide the richest multi-API
    material (HeCBench, Rodinia, RAJAPerf, BabelStream) implement most or all of
    them.

  \item \textbf{Portability periphery}: SYCL, Kokkos, RAJA, and OpenMP Target
    connect to the core cluster but with weaker edges. These APIs are
    concentrated in a smaller number of suites (primarily RAJAPerf, BabelStream,
    and HeCBench).

  \item \textbf{Isolated APIs}: OCCA, Thrust, Chapel, Coarray Fortran, Charm++,
    C++ AMP, UPC, OpenSHEM, HMPP, Python, HLS, GSParLib, STM/TM, and Serial
    appear as isolated or weakly connected nodes. Their co-occurrence with any
    other API falls below the edge threshold of~3.
\end{enumerate}

Figure~\ref{fig:api-coverage} presents the per-API coverage in a ranked bar
chart, making the quantitative dominance of the top-3 APIs clear.

\begin{figure}[htbp]
  \centering
  \includegraphics[width=\columnwidth]{api_coverage_bar_v5.png}
  \caption{API coverage across the expanded benchmark survey. OpenMP (${\sim}50$
    benchmarks), MPI (${\sim}45$), and CUDA (${\sim}38$) dominate, with OpenACC
    (${\sim}17$), HIP (${\sim}15$), and MPI+OpenMP (${\sim}15$) forming a second
    tier. OpenMP Target (${\sim}8$), SYCL (${\sim}7$), Kokkos (${\sim}7$), and
    RAJA (${\sim}7$) form a third tier. APIs below 5 benchmarks were excluded
    from evaluation consideration due to insufficient benchmark material for
    statistical evaluation.}
  \label{fig:api-coverage}
\end{figure}

\subsubsection*{Kernel-Level Co-Occurrence}

The repository-level analysis in Figures~\ref{fig:api-heatmap}--\ref{fig:api-coverage}
counts how many \emph{repositories} provide implementations in a given API pair,
but does not capture the depth of coverage within each repository. A suite like
HeCBench provides 516 individual kernels, each potentially implemented in
multiple APIs, while a single-kernel mini-application like miniBUDE provides
only~1 kernel across 12+ APIs. For translation evaluation, the kernel count is
the relevant unit: it determines how many independent evaluation instances are
available per direction.

\begin{figure}[htbp]
  \centering
  \includegraphics[width=\columnwidth]{kernel_level_cooccurrence_heatmap.png}
  \caption{Kernel-level API co-occurrence matrix. Each cell counts the number of
    individual kernels with verified equivalent implementations in both the row
    and column APIs. The CUDA--HIP pair leads with 633 shared kernels, followed
    by CUDA--SYCL (616), CUDA--OpenMP (472), and HIP--SYCL (615). OpenCL has
    dramatically lower kernel-level coverage: only 27 kernels share
    implementations with CUDA, and 24 with OpenMP. OpenMP Target provides 106
    shared kernels with each of CUDA, HIP, SYCL, and OpenMP, concentrated
    entirely in HeCBench. Sequential reference implementations exist for 106
    kernels.}
  \label{fig:kernel-cooccurrence}
\end{figure}

Figure~\ref{fig:kernel-cooccurrence} provides the strongest quantitative
evidence for API selection. The kernel-level data reveals a clear hierarchy:

\begin{itemize}
  \item \textbf{Tier~1 ($>$600 kernel pairs):} CUDA--HIP (633), CUDA--SYCL
    (616), HIP--SYCL (615). These pairs are dominated by HeCBench's systematic
    multi-API coverage.
  \item \textbf{Tier~2 (${\sim}450$ kernel pairs):} CUDA--OpenMP (472),
    HIP--OpenMP (453), SYCL--OpenMP (453). OpenMP's CPU threading model provides
    the deepest pool for \emph{paradigm-crossing} translation evaluation.
  \item \textbf{Tier~3 (${\sim}100$ kernel pairs):} OpenMP Target and Sequential
    at 106 each, concentrated in HeCBench.
  \item \textbf{Tier~4 ($<$30 kernel pairs):} OpenCL at 27 (CUDA--OpenCL) and 24
    (OpenMP--OpenCL), with an OpenCL diagonal of 27 total kernels.
\end{itemize}

The network visualization in Figure~\ref{fig:kernel-network} makes this tiered
structure visually apparent.

\begin{figure}[htbp]
  \centering
  \includegraphics[width=\columnwidth]{kernel_level_cooccurrence_network.png}
  \caption{Kernel-level API co-occurrence network. Edges shown for $\geq 20$
    kernel pairs; edge width proportional to count. The tight
    CUDA/HIP/SYCL/OpenMP cluster reflects HeCBench and RAJAPerf's uniform
    multi-API coverage. OpenCL appears as a peripheral node with weak connections
    (27 CUDA--OpenCL kernel pairs, just above the 20-kernel threshold). RAJA,
    stdpar, and Thrust appear as isolated nodes with fewer than 20 kernel-level
    pairs.}
  \label{fig:kernel-network}
\end{figure}

\subsection{Translation Difficulty Taxonomy}
\label{sec:appendix-a2}

The choice of API pairs for evaluation is motivated not only by data
availability (Section~\ref{sec:appendix-a1}) but by the \emph{structural
transformation complexity} each translation direction demands. We classify
translation directions into four categories:

\textbf{Syntactic renaming} (e.g., CUDA to HIP): Near-1:1 API call mapping.
\texttt{cudaMalloc} becomes \texttt{hipMalloc}; kernel launch syntax is
preserved verbatim; thread-index arithmetic (\texttt{threadIdx.x},
\texttt{blockIdx.x}) maps directly to HIP equivalents. This category tests API
name lookup, not parallel reasoning.

\begin{lstlisting}[language=C, caption={CUDA to HIP: syntactic renaming.}]
// CUDA                          // HIP (syntactic rename)
cudaMalloc(&d_a, size);          hipMalloc(&d_a, size);
kernel<<<grid, block>>>(d_a);    hipLaunchKernelGGL(kernel,
                                   grid, block, 0, 0, d_a);
cudaMemcpy(h_a, d_a, ...);      hipMemcpy(h_a, d_a, ...);
\end{lstlisting}

\textbf{Paradigm translation} (e.g., CUDA to OpenMP): SPMD to fork-join model
transformation. Explicit GPU thread indexing must be converted to loop-based
parallelism with OpenMP directives. Shared memory and synchronization primitives
(\texttt{\_\_shared\_\_}, \texttt{\_\_syncthreads()}) must be replaced with
stack-allocated arrays and implicit barrier semantics. Memory management
(\texttt{cudaMalloc}/\texttt{cudaMemcpy}/\texttt{cudaFree}) is entirely
eliminated.

\begin{lstlisting}[language=C, caption={CUDA to OpenMP: paradigm translation.}]
// CUDA                                  // OpenMP
__global__ void kern(float* a, int n) {  void kern(float* a, int n) {
  int i = blockIdx.x * blockDim.x         #pragma omp parallel for
           + threadIdx.x;                  for (int i = 0; i < n; i++) {
  if (i < n) a[i] = a[i] * 2.0f;            a[i] = a[i] * 2.0f;
}                                          }
                                         }
\end{lstlisting}

\textbf{Architecture split} (e.g., CUDA to OpenCL): Unified single-source CUDA
must be split into separate host code (OpenCL runtime API calls for
platform/device/context/queue management) and kernel source (\texttt{.cl} files
with OpenCL~C syntax). The memory model changes from CUDA's unified virtual
addressing to OpenCL's explicit buffer objects. Kernel launch syntax transforms
from \texttt{<<<grid, block>>>} to \texttt{clEnqueueNDRangeKernel} with explicit
work-group sizing.

\textbf{Directive insertion} (e.g., sequential~C to OpenMP): Adding parallelism
to sequential code by identifying parallelizable loops and inserting appropriate
\texttt{\#pragma omp} directives with correct data-sharing clauses
(\texttt{shared}, \texttt{private}, \texttt{reduction}). This is the inverse of
the paradigm translation direction and tests whether the LLM can \emph{discover}
parallelism rather than \emph{translate} it.

ParBench focuses on paradigm translation (CUDA/OpenMP/OpenCL
cross-translations) because this category maximizes the semantic challenge while
remaining tractable for automated verification. Syntactic renaming (CUDA to HIP)
is too easy to be informative; architecture split (CUDA to OpenCL) combines
paradigm translation with source-file restructuring, adding confounding
complexity.

\subsection{HIP and SYCL Exclusion Rationale}
\label{sec:appendix-a3}

Despite the large kernel-level co-occurrence counts for CUDA--HIP (633) and
CUDA--SYCL (616), these pairs were excluded from the primary evaluation for
distinct reasons.

\textbf{HIP exclusion:} CUDA-to-HIP translation is predominantly a syntactic
renaming task. The AMD HIP API was explicitly designed as a drop-in replacement
for CUDA, preserving the SPMD programming model, kernel launch syntax, and
memory management API. Automated tools such as \texttt{hipify-perl} and
\texttt{hipify-clang} achieve near-perfect conversion rates. Evaluating LLMs on
CUDA-to-HIP translation would measure API name lookup ability rather than
parallel programming reasoning---the core capability ParBench is designed to
assess.

\textbf{SYCL exclusion:} While SYCL introduces a genuinely different programming
model (single-source C++ with lambda-based kernel dispatch and buffer/accessor
memory management), it preserves the GPU execution model. CUDA-to-SYCL
translation is more complex than CUDA-to-HIP but remains within the same
architectural paradigm (GPU offload). Furthermore, SYCL benchmark coverage is
concentrated almost entirely in HeCBench, limiting the diversity of evaluation
instances. The CUDA-to-OpenMP direction provides a more fundamental paradigm
crossing (GPU SPMD to CPU fork-join) with better benchmark diversity across
Rodinia, HeCBench, XSBench, and RSBench.

\subsection{OpenACC Exclusion}
\label{sec:appendix-a4}

OpenACC was excluded despite appearing in approximately 17 benchmarks in the
expanded survey (Figure~\ref{fig:api-coverage}). Three factors motivated this
decision:

\begin{enumerate}
  \item \textbf{Limited co-occurrence with CUDA/OpenMP}: Of the 17 benchmarks
    with OpenACC, only~3 also provide both CUDA and OpenMP implementations. This
    yields insufficient material for statistically meaningful cross-API
    evaluation.

  \item \textbf{Paradigm overlap with OpenMP}: OpenACC's directive-based model
    (\texttt{\#pragma acc parallel loop}) occupies a similar conceptual niche to
    OpenMP's target offload directives (\texttt{\#pragma omp target teams
    distribute}). Including both would provide diminishing returns in
    programming-model diversity.

  \item \textbf{Compiler availability}: OpenACC compilation requires the NVIDIA
    HPC SDK (\texttt{nvc}) or GCC with \texttt{-fopenacc}. This is a less
    universally available toolchain than CUDA (\texttt{nvcc}) or OpenMP (any
    modern C/C++ compiler), introducing a confounding variable between compiler
    availability and LLM translation capability.
\end{enumerate}

\subsection{OpenMP Target as Case Study}
\label{sec:appendix-a5}

OpenMP target offload (\texttt{\#pragma omp target}) occupies a middle ground
between CPU OpenMP and GPU-native APIs. It uses the OpenMP directive model but
targets GPU execution, requiring the programmer to specify data mapping
(\texttt{\#pragma omp target data map(...)}) and work distribution
(\texttt{\#pragma omp teams distribute parallel for}) explicitly.

ParBench evaluates OpenMP target as a case study rather than a primary direction
for three reasons:

\begin{enumerate}
  \item \textbf{Compiler requirement}: OpenMP target compilation for NVIDIA GPUs
    requires the NVIDIA HPC compiler \texttt{nvc} (part of the NVIDIA HPC
    SDK~24.3), which is not universally available. Standard GCC/Clang OpenMP
    support targets only CPU threading. This introduces a confounding variable: a
    failed build may indicate compiler unavailability rather than translation
    quality.

  \item \textbf{Limited benchmark coverage}: OpenMP target implementations are
    available for XSBench, RSBench, and a subset of HeCBench curated kernels
    (106 kernel-level co-occurrences with CUDA, as shown in
    Figure~\ref{fig:kernel-cooccurrence}). Rodinia does not provide OpenMP target
    variants.

  \item \textbf{Compilation model difference}: OpenMP target generates GPU
    offload code through the compiler, while CPU OpenMP generates threaded CPU
    code. A translation from CUDA to OpenMP target preserves GPU execution
    semantics but changes the syntax entirely; a translation from CUDA to CPU
    OpenMP changes both the execution model and the syntax. These are
    qualitatively different evaluation targets that should not be conflated in
    aggregate statistics.
\end{enumerate}

The case study results for OpenMP target are reported separately in the
evaluation (Section~\ref{sec:results} of the main paper) to enable direct comparison with the
primary CUDA/OpenMP/OpenCL directions without confounding the aggregate pass
rates.


%% ============================================================================
%% APPENDIX B: Benchmark Kernel Survey
%% ============================================================================

\section{Benchmark Kernel Survey}
\label{sec:appendix-b}

This appendix documents the systematic survey and multi-stage selection process
that produced ParBench's evaluation corpus. The main paper (Section~\ref{sec:benchmark-curation})
summarizes this process; here we provide the complete data, selection criteria,
and exclusion rationale.

\subsection{Full Repository Inventory}
\label{sec:appendix-b1}

The survey identified 35 benchmark repositories spanning five categories of HPC
benchmark software. Figure~\ref{fig:benchmark-types} shows the distribution
across categories.

\begin{figure}[htbp]
  \centering
  \includegraphics[width=\columnwidth]{benchmark_types.png}
  \caption{Distribution of benchmark types across the 35 surveyed repositories.
    Suites (13, 37\%) and mini-applications (12, 34\%) dominate, together
    accounting for 71\% of the survey. Proxy applications (4, 11\%), libraries
    (4, 11\%), and full applications (2, 6\%) complete the distribution. Suites
    and mini-applications provide the richest kernel-level material for
    translation evaluation because they contain multiple independent
    computational kernels with self-checking verification.}
  \label{fig:benchmark-types}
\end{figure}

The relationship between benchmarks and APIs is not uniform: some repositories
implement a single API (e.g., SPEC OMP, PARSEC, LonestarGPU), while others
provide implementations across 10+ APIs (BabelStream with~12, miniBUDE
with~14). Figure~\ref{fig:bipartite} visualizes this heterogeneity as a
bipartite network.

% TODO: Restore or update benchmark_api_bipartite_color_readable_v7.png
% (deleted from analysis/visualizations/; exists in
%  analysis/visualizations/take in paper/benchmark_api_bipartite_color_readable_v7.png)
\begin{figure}[htbp]
  \centering
  % \includegraphics[width=\columnwidth]{benchmark_api_bipartite_color_readable_v7.png}
  \fbox{\parbox{0.9\columnwidth}{\centering\vspace{1em}\textbf{[Missing: benchmark\_api\_bipartite\_color\_readable\_v7.png]}\vspace{1em}}}
  \caption{Bipartite network linking benchmarks (circles, right) to parallel
    APIs (squares, left). Benchmark nodes are colored by type and sized
    proportional to the number of APIs they implement. API nodes are colored by
    family (vendor GPU, open standard GPU, directives, message passing, CPU
    parallel libs, portability libs, single-source C++, PGAS/one-sided). Green
    edges and green node outlines indicate verified download availability; grey
    indicates unverified. The dense connectivity between CUDA, OpenMP, HIP, and
    SYCL (left column) and the suite-type benchmarks (center) demonstrates the
    concentration of multi-API material in a small number of large suites,
    particularly HeCBench, RAJAPerf, and BabelStream.}
  \label{fig:bipartite}
\end{figure}

Figure~\ref{fig:api-count-hist} quantifies the API breadth distribution across
the surveyed benchmarks.

\begin{figure}[htbp]
  \centering
  \includegraphics[width=\columnwidth]{benchmark_api_count_hist_v5.png}
  \caption{Distribution of API breadth (number of parallel APIs supported)
    across surveyed benchmarks. The modal benchmark supports 3~APIs, with a peak
    of approximately 18 benchmarks in the 3-API bin. A long tail extends to 12+
    APIs, occupied by BabelStream and miniBUDE. Approximately 13 benchmarks
    support only 1--2 APIs (single-paradigm implementations). Benchmarks
    supporting 3 or more APIs are the primary candidates for multi-direction
    translation evaluation.}
  \label{fig:api-count-hist}
\end{figure}

The combination of Figures~\ref{fig:benchmark-types}--\ref{fig:api-count-hist}
establishes that the benchmark landscape is dominated by a small number of large
suites with broad API coverage. This Pareto distribution motivated the selection
strategy: focus on the richest suites (HeCBench, Rodinia) supplemented by
targeted mini-applications (XSBench, RSBench, mixbench) to ensure domain
diversity.

\subsection{Quality Tier Classification}
\label{sec:appendix-b2}

All 35 surveyed repositories were classified into two quality tiers based on
three criteria: build documentation, verification capability, and maintenance
status.

% TODO: Restore quality_tiers.png (deleted from analysis/visualizations/)
\begin{figure}[htbp]
  \centering
  % \includegraphics[width=\columnwidth]{quality_tiers.png}
  \fbox{\parbox{0.9\columnwidth}{\centering\vspace{1em}\textbf{[Missing: quality\_tiers.png]}\vspace{1em}}}
  \caption{Quality tier distribution across surveyed repositories. 30 of 35
    repositories (85\%) meet Tier~A criteria (documented build process, automated
    verification, and active maintenance within 2~years). 5 repositories (14\%)
    are classified as Tier~B due to partial verification, limited API coverage,
    or unmaintained status.}
  \label{fig:quality-tiers}
\end{figure}

\textbf{Tier~A criteria} (all three required):
\begin{itemize}
  \item Documented build process (Makefile, CMake, or equivalent with clear
    instructions)
  \item Automated verification: self-checking output (print PASS/FAIL, compute
    checksum, compare against reference) or reference output file comparison
  \item Active maintenance: commits within 2~years of the survey date, or stable
    release with no known build failures on modern toolchains
\end{itemize}

\textbf{Tier~B criteria} (any one sufficient for downgrade):
\begin{itemize}
  \item Partial verification only (e.g., visual output inspection, manual
    comparison)
  \item Limited API coverage (single API, reducing cross-translation utility)
  \item Unmaintained but still buildable (no commits in 2+ years, but Makefiles
    still produce working binaries)
\end{itemize}

Tier~B repositories were not excluded from the survey but were deprioritized in
kernel selection. Their kernels were considered only when no Tier~A alternative
existed for a given computational domain.

\subsection{Candidate Ranking and HeCBench Selection Funnel}
\label{sec:appendix-b3}

The selection of individual kernels for ParBench's evaluation corpus required
balancing two competing objectives: \emph{kernel richness} (maximizing the
number of independent evaluation instances) and \emph{API breadth} (maximizing
the number of translation directions per kernel). Figure~\ref{fig:kernel-api-bubble}
plots these two dimensions for all surveyed repositories.

\begin{figure}[htbp]
  \centering
  \includegraphics[width=\columnwidth]{kernel_api_bubble.png}
  \caption{Benchmark kernel richness (y-axis, capped at 100 for readability) vs.\
    API breadth (x-axis, number of APIs supported). Bubble size proportional to
    total kernel count; color indicates benchmark type. HeCBench (516 kernels,
    4~APIs) dominates the upper region. RAJAPerf (80 kernels, 5~APIs) is the
    second-largest. CloverLeaf (23 kernels, 9~APIs) and BabelStream (5 kernels,
    12~APIs) occupy the high-API-breadth region but with far fewer kernels. The
    optimal candidates for corpus construction lie in the high-kernel,
    moderate-API quadrant (upper-center).}
  \label{fig:kernel-api-bubble}
\end{figure}

Figure~\ref{fig:top-candidates} presents the top candidates ranked by their
combined kernel-count and API-breadth scores, revealing the fundamental tradeoff
in corpus construction.

\begin{figure}[htbp]
  \centering
  \includegraphics[width=\columnwidth]{top_candidates.png}
  \caption{Top ``Rosetta Stone'' candidates ranked by kernel count (blue bars,
    scaled for readability; true counts annotated above) and API breadth (red
    bars). HeCBench provides the largest kernel pool (516 kernels) but with
    only~4 primary APIs (CUDA, HIP, SYCL, OpenMP). BabelStream provides the
    broadest API coverage (12~APIs) but with only 5~kernels. RAJAPerf (80
    kernels, 7~APIs) offers a balanced middle ground. Rodinia (24 kernels, 3
    primary APIs) provides established, well-understood benchmarks with strong
    community recognition.}
  \label{fig:top-candidates}
\end{figure}

\subsubsection*{HeCBench Selection Funnel}

HeCBench's 516 kernels were filtered through a multi-stage funnel to identify
the curated evaluation set:

\begin{itemize}
  \item \textbf{Stage~1 (API filter):} 516 total kernels in the HeCBench
    repository. Of these, 272 provide implementations in all 4~primary APIs
    (CUDA, HIP, SYCL, OpenMP) with Makefiles and self-checking output. The
    remaining 244 lack one or more API variants, Makefiles in the CUDA variant,
    or self-checking output patterns.

  \item \textbf{Stage~2 (Complexity filter):} The 272 candidates were filtered
    for evaluation tractability. Kernels with more than 15 source files or fewer
    than 2 were excluded (the former are too complex for single-prompt
    translation; the latter are trivial boilerplate). Kernels requiring external
    input data files were also excluded to ensure self-contained evaluation.

  \item \textbf{Stage~3 (Domain diversity):} From the complexity-filtered set,
    60 kernels were selected to maximize domain diversity across computational
    categories: linear algebra, stencil computation, graph algorithms, machine
    learning, physics simulation, sorting, reduction, scan, image processing,
    and cryptography.

  \item \textbf{Stage~4 (Curation):} The 60-kernel working set was narrowed to a
    final curated set of 20 kernels based on manual inspection of build
    reliability, output determinism, and domain representativeness. Of these, 10
    kernels were included in the current ParBench evaluation corpus, with the
    remaining~10 reserved for future expansion.
\end{itemize}

The final 20 curated HeCBench kernels span diverse computational domains and
algorithmic patterns, ensuring that the evaluation corpus tests LLM translation
capability across a representative sample of real HPC workloads.

\subsection{Rodinia Kernel Inventory}
\label{sec:appendix-b4}

Rodinia provides the foundation of ParBench's evaluation corpus. As one of the
most widely cited HPC benchmark suites (originally published at IISWC
2009~\cite{Rodinia2009}), it offers strong community recognition and
well-understood computational characteristics. ParBench uses 22 Rodinia kernels
across three APIs (CUDA, OpenMP, OpenCL), yielding 60 total specs (some kernels
lack one or more API variants).

Of these 60 specs:
\begin{itemize}
  \item \textbf{54 TRUE PASS}: Verified correct output with exit code~0 and
    stdout pattern matching using conjunction verification (both conditions must
    hold).
  \item \textbf{6 KNOWN\_FAIL}: Excluded from evaluation batches due to
    pre-existing build or runtime failures unrelated to LLM translation. These
    include 2~kernels using deprecated CUDA \texttt{texture<>} references
    removed in CUDA~12 (\texttt{kmeans-cuda}, \texttt{mummergpu-cuda}),
    1~kernel requiring OpenGL dependencies (\texttt{hybridsort-cuda}), 1~kernel
    with CUDA API dependencies in its OpenMP variant (\texttt{mummergpu-omp}),
    and 2~kernels with pre-existing OpenCL runtime issues (\texttt{nn-opencl},
    \texttt{kmeans-opencl}).
\end{itemize}

\subsection{XSBench, RSBench, and mixbench Profiles}
\label{sec:appendix-b5}

Three additional benchmark suites complement Rodinia and HeCBench in the
evaluation corpus:

\textbf{XSBench} (4 specs: CUDA, OpenMP, OpenCL, OpenMP Target): A Monte Carlo
neutron cross-section lookup proxy application from Argonne National Laboratory.
XSBench implements a hash-based unionized energy grid algorithm that is
representative of nuclear reactor simulation workloads. All 4 API variants pass
baseline verification. The OpenMP variant produces a history-based checksum
(941535); the CUDA, OpenCL, and OpenMP Target variants produce an event-based
checksum (945990). This asymmetry reflects algorithmic differences between the
CPU and GPU implementations rather than correctness issues.

\textbf{RSBench} (4 specs: CUDA, OpenMP, OpenCL, OpenMP Target): A simplified
variant of XSBench that isolates the multipole cross-section lookup kernel.
RSBench was added to ParBench to enable controlled comparison with XSBench on
the same algorithmic family at different complexity levels.

\textbf{mixbench} (3 specs: CUDA, OpenMP, OpenCL): A micro-benchmark designed to
measure GPU arithmetic throughput at varying compute-to-memory ratios. mixbench
exercises the roofline boundary, making it valuable for evaluating whether LLM
translations preserve performance-critical arithmetic intensity.

\subsection{Exclusion Log and Verification Methods}
\label{sec:appendix-b6}

Figure~\ref{fig:verification-network} groups the surveyed benchmarks by their
verification approach, illustrating the diversity of correctness checking
methods across the survey and the basis for exclusion decisions.

\begin{figure}[htbp]
  \centering
  \includegraphics[width=\columnwidth]{verification_network.png}
  \caption{Benchmarks grouped by verification method. Large colored nodes
    represent verification approaches: checksum comparison, convergence testing,
    reference output file comparison, output-to-standard comparison, and
    self-validation. Benchmark nodes (smaller, colored by type) are connected to
    their verification method(s). Benchmarks without automated verification
    appear as weakly connected or isolated nodes and were excluded or downgraded
    to Tier~B. The checksum and reference output clusters contain the majority
    of Tier~A benchmarks.}
  \label{fig:verification-network}
\end{figure}

\subsubsection*{Repository-Level Exclusions}

Five of the initial 40 candidate repositories were excluded during the survey
phase:
\begin{itemize}
  \item \textbf{Download failures} (2): Repository URLs were dead or
    access-restricted at the time of survey.
  \item \textbf{Insufficient documentation} (2): No build instructions, no
    Makefiles, or build system required proprietary dependencies not available on
    the evaluation platform.
  \item \textbf{Duplicate content} (1): One repository was a fork of an
    already-surveyed suite with no additional API variants.
\end{itemize}

\subsubsection*{Kernel-Level Exclusions}

Within the selected suites, individual kernels were excluded for the following
reasons:
\begin{itemize}
  \item \textbf{Missing API variants}: Kernels available in only one API (no
    cross-translation possible).
  \item \textbf{No self-checking output}: Kernels that produce graphical output,
    require manual inspection, or have no deterministic expected output.
  \item \textbf{Excessive complexity}: Kernels requiring more than 15 source
    files, external data files not included in the repository, or runtime
    dependencies (e.g., OpenGL, MPI multi-node) not available on the single-GPU
    evaluation platform.
  \item \textbf{Non-deterministic output}: Kernels whose output depends on
    execution order, random seeds without fixed initialization, or
    floating-point non-associativity beyond the tolerance of the verification
    method.
\end{itemize}

The complete exclusion log with per-kernel justifications is maintained in
ParBench's specification metadata (\texttt{specs/*.json}) and the project's
known-issues documentation.


%% ============================================================================
%% APPENDIX C: Detailed Evaluation Findings
%% ============================================================================

\section{Detailed Evaluation Findings}
\label{sec:appendix-c}

This appendix provides per-kernel and per-direction breakdowns of the evaluation
results summarized in Section~\ref{sec:results} of the main paper, including self-repair
dynamics, augmentation transform statistics, the HeCBench selection funnel, and
a detailed case study of multi-file translation consistency failures.

\subsection{Self-Repair Transition Matrices}
\label{sec:appendix-c1}

ParBench's multi-attempt evaluation protocol allows LLMs to receive
compiler/runtime error feedback and attempt corrections. Figure~\ref{fig:repair-matrix}
presents the self-repair transition matrices, showing how initial failure modes
map to final evaluation outcomes across all models and directions.

\begin{figure}[htbp]
  \centering
  \includegraphics[width=\columnwidth]{figures/c1_repair_transition_matrix.pdf}
  \caption{Self-repair transition matrix showing initial failure mode to final
    outcome counts across all translation directions. Rows represent the failure
    mode of the first attempt (BUILD\_FAIL, RUN\_FAIL, VERIFY\_FAIL); columns
    represent the final evaluation outcome after up to 3~attempts. Off-diagonal
    entries indicate successful self-repair transitions.}
  \label{fig:repair-matrix}
\end{figure}

The transition matrices reveal several key dynamics of LLM self-repair behavior:

\begin{itemize}
  \item \textbf{Initial failure mode distribution}: The majority of first-attempt
    failures are BUILD\_FAIL (missing headers, wrong API calls, syntax errors),
    followed by RUN\_FAIL (segmentation faults, timeouts) and VERIFY\_FAIL
    (incorrect output). This ordering reflects the pipeline's sequential
    verification: build errors are caught before runtime errors, which are caught
    before output verification.

  \item \textbf{Repair rate asymmetry}: BUILD\_FAIL has the highest repair rate,
    as compiler error messages provide actionable feedback (missing header name,
    undeclared identifier, type mismatch). RUN\_FAIL has a moderate repair rate
    when error messages are informative (e.g., ``segmentation fault at line~X'')
    but low repair rate for opaque failures (e.g., timeouts with no diagnostic
    output). VERIFY\_FAIL has the lowest repair rate, as output mismatches
    provide minimal diagnostic signal---the model receives ``output did not
    match expected pattern'' but lacks information about \emph{which} computation
    diverged.

  \item \textbf{Token consumption}: Self-repair attempts consume additional
    tokens for both the error feedback prompt and the corrected translation.
    Evaluations requiring 3~attempts consume approximately $2.5{\times}$ the
    tokens of first-attempt successes. However, the marginal cost of self-repair
    is low relative to the cost of a completely new translation, as the model can
    build on its previous attempt.

  \item \textbf{Attempt convergence}: Among evaluations that eventually pass,
    the median number of attempts is~2, indicating that most repairable failures
    are resolved on the second try. Third-attempt successes are rare, suggesting
    diminishing returns from additional retry budget beyond~3.
\end{itemize}

\subsection{Self-Repair Rate by Direction}
\label{sec:appendix-c2}

\begin{figure}[htbp]
  \centering
  \includegraphics[width=\columnwidth]{figures/c2_repair_rate_by_direction.pdf}
  \caption{Self-repair rate by translation direction with attempt count
    distribution. Each bar shows the fraction of initially-failing evaluations
    that are rescued by subsequent attempts, broken down by the number of
    attempts required. Directions are ordered by repair rate.}
  \label{fig:repair-direction}
\end{figure}

The repair rate varies significantly by translation direction
(Figure~\ref{fig:repair-direction}). Paradigm-preserving directions
(CUDA-to-OpenMP, OpenMP-to-CUDA) show higher repair rates than cross-paradigm
directions involving OpenCL, consistent with the structural similarity of CUDA
and OpenMP parallel patterns despite their different execution models.

This directional asymmetry in repair effectiveness has practical implications
for LLM-based translation tooling. A self-repair budget of 3~attempts is
sufficient to capture most recoverable failures in CUDA$\leftrightarrow$OpenMP
directions, but provides diminishing returns for OpenCL targets where the
failure modes (missing address space qualifiers, C++ to C99 decomposition) are
structurally harder to diagnose from compiler error messages alone.

Preliminary pass@$k$ experiments further illuminate the self-repair mechanism.
Comparing the primary campaign (temperature$\,{=}\,0.0$, 3~retries with
compiler feedback) against stochastic sampling (temperature$\,{=}\,0.7$,
1~attempt, 3~independent samples) on CUDA-to-OpenMP reveals a $2.3{\times}$
improvement from self-repair with greedy decoding (63.6\% vs.\ 27.3\% pass
rate). Notably, 7~of~11 kernel triplets produce identical outcomes across all
3~stochastic samples, indicating that sampling diversity alone is insufficient
to overcome systematic translation errors. Self-repair with compiler error
feedback is the dominant success factor, not stochastic exploration of the
output distribution.

\subsection{Augmentation Transform Frequency Per Kernel}
\label{sec:appendix-c3}

ParBench's code augmentation pipeline applies six AST-level transforms at four
intensity levels (L1--L4). The transforms do not apply uniformly: some kernels
have many candidate sites for a given transform, while others have few or none.
Figure~\ref{fig:transform-freq} presents the per-kernel, per-transform frequency
heatmap.

\begin{figure}[htbp]
  \centering
  \includegraphics[width=\columnwidth]{figures/c3_transform_frequency.pdf}
  \caption{Transform frequency heatmap showing candidate sites per kernel per
    transform type. Rows represent kernels; columns represent the six AST-level
    transforms. Cell intensity indicates the number of candidate sites
    identified by the augmentation pipeline. The distribution is highly
    non-uniform: control-flow-heavy kernels (e.g., \texttt{bfs},
    \texttt{particlefilter}) have many SwapCondition sites, while
    compute-intensive kernels (e.g., \texttt{hotspot}, \texttt{lud}) are
    dominated by ArithmeticTransform sites.}
  \label{fig:transform-freq}
\end{figure}

Aggregate transform frequencies across the corpus reveal a clear dominance
hierarchy:

\begin{itemize}
  \item \textbf{SwapCondition}: 162 candidate sites (most common). Targets
    \texttt{if/else} branches and ternary expressions. Its prevalence reflects
    the ubiquity of conditional logic in HPC kernels---boundary checks, loop
    guards, and branch-based algorithm selection all contribute candidate sites.

  \item \textbf{ArithmeticTransform}: 69 candidate sites. Applies commutative
    rewriting of arithmetic expressions (e.g., $a + b \rightarrow b + a$,
    $a \times (b + c) \rightarrow a \times b + a \times c$). Less common than
    SwapCondition because many arithmetic expressions in HPC code involve
    non-commutative operations (subtraction, division) or are too simple
    (single-operand) to transform meaningfully.

  \item \textbf{ChangeNames}: 55 candidate sites. Performs variable and
    parameter renaming while preserving semantic equivalence. The count is
    bounded by the number of local variables and parameters per kernel, which is
    typically 10--30 for Rodinia kernels.

  \item \textbf{TypedefExpansion}: 7 candidate sites. Inlines
    \texttt{typedef} aliases to their underlying types. Rare because most
    Rodinia kernels use primitive types directly rather than defining aliases.

  \item \textbf{PointerArithmeticToArrayIndex}: 6 candidate sites. Rewrites
    \texttt{*(ptr + i)} to \texttt{ptr[i]} and vice versa. Limited by the
    prevalence of pointer arithmetic in the corpus---most kernels use array
    indexing by default.

  \item \textbf{ChangeFunctionNames}: 2 candidate sites (least common). Limited
    by the function definition count in single-kernel files, which typically
    contain only 1--3 user-defined functions.
\end{itemize}

The level-invariance property (54/60 Rodinia specs pass at all levels L1--L4)
confirms that augmentation introduces no new correctness failures. This is
expected: all transforms preserve semantic equivalence by construction, and the
augmentation unit tests (\texttt{c\_augmentation/test\_transforms.py}, 15~tests)
verify this property for each transform in isolation. From an evaluation
perspective, augmentation levels serve as a robustness probe---demonstrating
that LLM translation quality is stable under semantics-preserving source
perturbations---rather than as independent samples that inflate the effective
sample size.

\subsection{HeCBench Kernel Selection Funnel}
\label{sec:appendix-c4}

\begin{figure}[htbp]
  \centering
  \includegraphics[width=\columnwidth]{figures/c4_selection_funnel.pdf}
  \caption{HeCBench kernel selection funnel. Starting from 513 kernels spanning
    CUDA, HIP, SYCL, and OpenMP, successive filters for API completeness (327
    with all 4~variants), build infrastructure (325), self-checking output
    patterns (242), complexity bounds, and domain diversity yield the 60-kernel
    working set and final 20-kernel curated evaluation corpus.}
  \label{fig:selection-funnel}
\end{figure}

HeCBench is the largest source of multi-API kernel implementations in the
benchmark landscape survey, with 513 kernels spanning CUDA, HIP, SYCL, and
OpenMP (Figure~\ref{fig:selection-funnel}). Its scale necessitated a structured
selection funnel to identify kernels suitable for automated evaluation:

\begin{enumerate}
  \item \textbf{API completeness} (513 $\rightarrow$ 327): Retained only kernels
    providing implementations in all 4~primary APIs. Kernels missing one or more
    API variants cannot support the full matrix of cross-API translation
    directions that ParBench evaluates.

  \item \textbf{Build infrastructure} (327 $\rightarrow$ 325): Excluded 2~kernels
    whose CUDA variants lacked Makefiles or had broken build configurations.
    All retained kernels build successfully with \texttt{make} on the evaluation
    platform (NVIDIA HPC SDK~24.3, CUDA~12.3).

  \item \textbf{Self-checking output} (325 $\rightarrow$ 242): Excluded kernels
    without deterministic, machine-parseable output. Retained kernels print
    checksums, PASS/FAIL indicators, or numerical results that can be matched
    against regular expression patterns for automated verification.

  \item \textbf{Complexity bounds} (242 $\rightarrow$ 60): Excluded kernels with
    more than 15 source files (too complex for single-prompt LLM translation) or
    fewer than 2~source files (trivial boilerplate). Also excluded kernels
    requiring external input data files not bundled with the repository.

  \item \textbf{Domain diversity and curation} (60 $\rightarrow$ 20): Manual
    inspection for build reliability, output determinism, and domain
    representativeness. The final 20~kernels span linear algebra, stencil
    computation, graph algorithms, machine learning, physics simulation, sorting,
    reduction, scan, image processing, and cryptography. Ten kernels are included
    in the current evaluation corpus; ten are reserved for future expansion.
\end{enumerate}

This funnel illustrates a general challenge in constructing LLM code translation
benchmarks: the vast majority of available HPC kernels are unsuitable for
automated evaluation due to missing API variants, non-deterministic output, or
excessive complexity. Only 3.9\% of HeCBench's 513 kernels (20/513) survive the
full selection pipeline, highlighting the gap between ``available parallel code''
and ``evaluable parallel code.''

\subsection{Case Study: Multi-File Translation Consistency Failure}
\label{sec:appendix-c5}

Full-program translations---where the LLM must rewrite both kernel code and host
code---expose a class of errors absent from kernel-only translations. We
document this failure class through a detailed case study of the
\texttt{rodinia-lavamd} kernel translated from OpenCL to CUDA.

\subsubsection*{Observation}

Across all 5~augmentation levels (L0--L4), the model (Qwen~3.5 397B) translates
OpenCL kernel computational patterns to CUDA correctly (\texttt{\_\_global\_\_},
\texttt{blockIdx.x}, \texttt{threadIdx.x}, \texttt{\_\_syncthreads()}) but
\textbf{fails to produce a compilable multi-file CUDA program} in all 15
attempts (5~levels $\times$ 3~retries). The failure modes vary across
augmentation levels---this is not a single consistent bug, but a family of
related multi-file consistency errors.

\subsubsection*{Per-Level Error Taxonomy}

\begin{table}[htbp]
  \centering
  \caption{Multi-file translation error taxonomy for \texttt{lavamd}
    OpenCL-to-CUDA across augmentation levels. The failure mode varies by level,
    demonstrating that augmentation perturbs \emph{which} errors the model makes
    without changing the outcome (BUILD\_FAIL at all levels).}
  \label{tab:lavamd-errors}
  \small
  \begin{tabular}{@{}clp{4.2cm}@{}}
    \toprule
    Level & Naming Correct? & Primary Failure Mode \\
    \midrule
    L0 & No  & Identifier renaming failure (retained \texttt{kernel\_gpu\_opencl}) \\
    L1 & Yes & Phantom dependencies (\texttt{timing.h}, \texttt{checkCUDAError}) \\
    L2 & No  & Identifier renaming failure (retained source name) \\
    L3 & Yes & Undefined helpers (\texttt{setdevice}, \texttt{checkCUDAError}) \\
    L4 & Yes & Undefined helpers + missing kernel launch syntax \\
    \bottomrule
  \end{tabular}
\end{table}

Table~\ref{tab:lavamd-errors} reveals that the ``semantic renaming failure''
(keeping \texttt{kernel\_gpu\_opencl} in CUDA code) only occurs at L0 and L2. At
L1, L3, and L4, the model correctly renames the function but fails on different
issues: phantom include paths, undefined helper functions, and missing
\texttt{<<<>>>} kernel launch syntax.

\subsubsection*{Retry Behavior}

All 5~levels share the same first-attempt error: the model hallucinates a
non-existent header (\texttt{timing.h}). Subsequent attempts diverge by level,
introducing \emph{new} error categories without converging toward a fix. This
non-convergent retry pattern---where each attempt produces qualitatively
different errors rather than incrementally correcting the previous
attempt---demonstrates a fundamental limitation of compiler-error-feedback-based
self-repair for multi-file consistency errors. The error messages are locally
informative (``undefined reference to \texttt{checkCUDAError}'') but do not
convey the systemic issue: the model must generate a coherent program scaffold,
not just fix individual compilation errors.

\subsubsection*{Implications}

This case study identifies four distinct error subcategories within the
``multi-file translation consistency failure'' class:

\begin{enumerate}
  \item \textbf{Phantom dependencies}: The model hallucinates include files
    (\texttt{timing.h}) or helper functions (\texttt{checkCUDAError},
    \texttt{setdevice}) that exist in the model's training data but not in the
    target program's file structure.

  \item \textbf{Cross-file identifier inconsistency}: The kernel file uses the
    target API's naming convention (\texttt{kernel\_gpu\_cuda}) while the wrapper
    file retains the source API's name (\texttt{kernel\_gpu\_opencl}), causing
    linker errors.

  \item \textbf{Missing host-side launch syntax}: The wrapper calls the kernel
    as a regular C function (\texttt{kernel\_gpu\_cuda(...)}) rather than using
    CUDA kernel launch syntax (\texttt{kernel\_gpu\_cuda<<<grid, block>>>(...)}).

  \item \textbf{Scaffold generation failure}: The model cannot synthesize the
    boilerplate infrastructure (device selection, error checking, memory
    allocation/deallocation) needed for a complete CUDA program, instead relying
    on non-existent helper functions.
\end{enumerate}

These findings suggest that a hybrid pipeline---combining LLM translation for
computational kernels with template-based scaffold generation for host
infrastructure---may outperform end-to-end LLM translation for multi-file
programs. However, the varying failure modes across augmentation levels indicate
that a simple post-processing pass (e.g., find-and-replace renaming) would only
address subcategory~2, leaving subcategories 1, 3, and~4 unresolved.

%% ============================================================================
%% APPENDIX D: Supplementary Tables and Figures
%% ============================================================================

\section{Supplementary Tables and Figures}
\label{sec:appendix-d}

This appendix contains detailed evidence tables and figures referenced from the main text.
Floats are ordered by their original section appearance.

% --- From Section 3 (ParBench Framework) ---

% src: tab:augmentation-levels verified against LEVEL_FRACTIONS = {1: 0.0, 2: 0.33, 3: 0.66, 4: 1.0} (01-01 verified 2026-04-03)
% src: in c_augmentation/augment_dataset.py line 111 (LEVEL_FRACTIONS dict)
% src: Transform selection (L1: 1 of 6, L2: max(1,int(6*0.33))=1, L3: max(1,int(6*0.66))=3, L4: all 6) verified against augment_code() lines 1796-1803
% src: Candidate fraction within each transform verified against _select_fraction() lines 152-167: max(1, floor(n*fraction))

\begin{table}[htbp]
\centering
\caption{Augmentation level definitions.}
\label{tab:augmentation-levels}
\small
\begin{tabular}{@{}cccp{3.8cm}@{}}
\toprule
Level & Transform selection & Candidate fraction & Description \\
\midrule
L0 & None (original) & 0\% & Unmodified source code (baseline) \\
L1 & 1 transform & 1 site per transform & Minimal surface perturbation \\
L2 & 33\% of transforms & 33\% of candidates & Moderate obfuscation \\
L3 & 66\% of transforms & 66\% of candidates & Heavy obfuscation \\
L4 & All transforms & 100\% of candidates & Maximum obfuscation \\
\bottomrule
\end{tabular}
\end{table}

% --- From Section 4 (Benchmark Curation) ---


\begin{figure}[htbp]
\centering
\includegraphics[width=\columnwidth]{f2_repo_vs_kernel.pdf}
\caption{Repository-level vs.\ kernel-level translation pair counts. Left bars show the number of repositories containing both APIs; right bars show the number of independent kernel-level translation pairs. The multipliers (79$\times$--308$\times$) demonstrate that repository-level counting substantially underestimates the translation evaluation opportunity.}
\label{fig:repo-vs-kernel}
\end{figure}


\begin{table}[htbp]
\centering
\caption{Survey---kernel-level translation pair counts.}
\label{tab:kernel-pairs}
\small
\begin{tabular}{@{}lccl@{}}
\toprule
API Pair & Repos & Kernel Pairs & Primary Sources \\
\midrule
CUDA--HIP & 3 & 633 & HeCBench (504), RAJAPerf (106), CloverLeaf (16) \\
CUDA--SYCL & 2 & 616 & HeCBench (487), RAJAPerf (106), CloverLeaf (16) \\
CUDA--OpenMP & 6 & 472 & HeCBench (325), RAJAPerf (106), Rodinia (19) \\
HIP--OpenMP & 2 & 453 & HeCBench (324), RAJAPerf (106), CloverLeaf (16) \\
CUDA--OpenCL & 6 & ${\sim}$200 & Rodinia (19), Parboil, SHOC \\
\bottomrule
\end{tabular}
\end{table}


\begin{table*}[htbp]
\centering
\caption{Programming model characteristics of the three primary translation APIs. These differences define the structural transformation challenges that LLM-based translation must address.}
\label{tab:api-characteristics}
\small
\begin{tabular}{@{}lp{4.0cm}p{4.0cm}p{4.0cm}@{}}
\toprule
Property & CUDA & OpenMP & OpenCL \\
\midrule
Execution model & SIMT (warps of 32 threads) & Fork-join thread teams & NDRange work-groups \\
Memory management & Explicit device allocation (\texttt{cudaMalloc}/\texttt{cudaMemcpy}) & Implicit shared CPU memory & Explicit buffer objects (\texttt{clCreateBuffer}/\texttt{clEnqueueReadBuffer}) \\
Kernel specification & Compiled with host code (\texttt{\_\_global\_\_} functions) & Pragma directives (\texttt{\#pragma omp parallel for}) & Runtime string compilation (\texttt{clCreateProgramWithSource}) \\
Source structure & Single or multi-file (51\% multi-file in corpus) & Single file (directives annotate sequential code) & Separate host program + \texttt{.cl} kernel files \\
Thread indexing & \texttt{threadIdx.x + blockIdx.x * blockDim.x} & Loop iteration variable or \texttt{omp\_get\_thread\_num()} & \texttt{get\_global\_id(0)} \\
\bottomrule
\end{tabular}
\end{table*}


\begin{table}[htbp]
\centering
\caption{Suite summary showing total specs, verified PASS, KNOWN\_FAIL, and API coverage for each of the five benchmark suites.}
\label{tab:suite-summary}
\small
% src: Verified 2026-04-03 against specs/ on disk:
% src: Rodinia: ls specs/rodinia-*.json | wc -l => 60; unique kernels => 22; KF=6 (known-issues.md)
% src: HeCBench: 10 curated kernels x {cuda,omp,omp_target} = 25 specs (from 135 total); KF=2
% src: XSBench: 4; RSBench: 4; mixbench: 3 (all verified via ls specs/{suite}-*.json)
% src: Sums: 22+10+1+1+1=35 kernels, 60+25+4+4+3=96 specs, 54+23+4+4+3=88 PASS, 6+2+0+0+0=8 KF
\begin{tabular}{@{}lccccc@{}}
\toprule
Suite & Kernels & Total & PASS & KF & APIs \\
\midrule
Rodinia & 22 & 60 & 54 & 6 & CUDA, OMP, OCL \\
HeCBench (curated) & 10 & 25 & 23 & 2 & CUDA, OMP, OMP-tgt \\
XSBench & 1 & 4 & 4 & 0 & CUDA, OMP, OCL, OMP-tgt \\
RSBench & 1 & 4 & 4 & 0 & CUDA, OMP, OCL, OMP-tgt \\
mixbench & 1 & 3 & 3 & 0 & CUDA, OMP, OCL \\
\midrule
\textbf{Total} & \textbf{35} & \textbf{96} & \textbf{88} & \textbf{8} & \\
\bottomrule
\end{tabular}
\end{table}


\begin{table}[htbp]
\centering
\caption{Quantitative characterization of the evaluation corpus. SLoC counts physical source lines in the CUDA prompt payload. Multi-File gives specs requiring translation of $>$1 source file. Std.\ shows the dominant language standard per suite.}
\label{tab:benchmark-characterization}
\small
% src: Rodinia SLoC 195--3304, median 333.5 (rounded to 334), 21/60 multi-file; verified 2026-04-03
% src: HeCBench SLoC 80--235, median 180, 0/25 multi-file; verified 2026-04-03
% src: XSBench 1390, RSBench 1016, mixbench 312 SLoC; verified 2026-04-03
\begin{tabular}{@{}lrcrcl@{}}
\toprule
Suite & Kernels & SLoC Range & Med. & Multi-File & Std. \\
\midrule
Rodinia         & 22 & 195--3{,}304 & 334 & 21/60 (35\%) & C++14 \\
HeCBench (curated) & 10 &    80--235    & 180 &  0/25 ~(0\%) & C++17 \\
XSBench         &  1 &    1{,}390    &  -- &  2/4~~(50\%) & C11   \\
RSBench         &  1 &    1{,}016    &  -- &  1/4~~(25\%) & C11   \\
mixbench        &  1 &       312     &  -- &  0/3~~~(0\%) & C++11 \\
\midrule
\textbf{Total}  & \textbf{35} & \textbf{80--3{,}304} & \textbf{271} & \textbf{24/96 (25\%)} & \\
\bottomrule
\end{tabular}
\end{table}


\begin{table}[htbp]
\centering
\caption{Domain category distribution across the 35-kernel evaluation corpus. Categories are assigned from the canonical \texttt{sloc\_analysis.json} classification. Suite annotations show which suites contribute kernels to each category.}
\label{tab:category-distribution}
\small
\begin{tabular}{@{}lrl@{}}
\toprule
Category & Kernels & Contributing Suites \\
\midrule
Physics simulation     & 7 & Rodinia (5), XSBench, RSBench \\
Graph algorithms       & 5 & Rodinia (3), HeCBench (2) \\
Stencil computation    & 5 & HeCBench (5) \\
Image processing       & 2 & Rodinia (2) \\
Linear algebra         & 2 & Rodinia (2) \\
Machine learning       & 2 & Rodinia (2) \\
Molecular dynamics     & 2 & Rodinia (1), HeCBench (1) \\
Reduction              & 1 & HeCBench (1) \\
Sort                   & 1 & Rodinia (1) \\
Other                  & 8 & Rodinia (6), HeCBench (1), mixbench (1) \\
\midrule
\textbf{Total}         & \textbf{35} & \textbf{5 suites} \\
\bottomrule
\end{tabular}
\end{table}

% --- From Section 5 (Experimental Setup) ---


% src: 2 models only: Qwen 3.5 397B-A17B (Together AI) + GPT-4.1 mini (Azure OpenAI); verified 2026-04-03
% src: GPT-4.1 mini column has \pending markers (Le's data pending); verified 2026-04-03
\begin{table}[htbp]
\centering
\caption{Model configurations.}
\label{tab:model-config}
\small
\begin{tabular}{@{}llllc@{}}
\toprule
Model & Provider & Arch. & Parameters & Temp. \\
\midrule
Qwen 3.5 397B-A17B & Together AI & MoE & 397B / 17B active & 0 \\
GPT-4.1 mini & Azure OpenAI & N/D & \pending{size} & 0 \\
\bottomrule
\end{tabular}
\end{table}


\begin{table}[htbp]
\centering
\caption{Hardware and software configuration.}
\label{tab:hardware}
\small
\begin{tabular}{@{}lll@{}}
\toprule
Component & Qwen (primary) & GPT-4.1 mini (collaborator) \\
\midrule
GPU & RTX 4070 (12 GB) & \pending{GPU} \\
CPU & Ryzen 9 7900X & \pending{CPU} \\
OS & Ubuntu 24.04 LTS & \pending{OS} \\
CUDA toolkit & HPC SDK 24.3 & \pending{toolkit} \\
C/C++ compiler & GCC 12.4 & \pending{compiler} \\
OpenCL runtime & NVIDIA HPC SDK & \pending{OCL runtime} \\
Python & 3.12.3 & \pending{Python} \\
\bottomrule
\end{tabular}
\end{table}

% --- From Section 6 (Results) ---


% src: paper_data.json > primary_campaign > self_repair > per_initial_failure (all rows verified 2026-04-05)
% src: total repaired 72+17+8+15+0=112; total initial failures 346+148+28+27+1=550; 550=710-160 verified
\begin{table}[htbp]
\centering
\caption{Self-repair failure mode transitions (all tasks, Qwen~3.5; GPT-4.1~mini pending). Rows: initial failure mode. Columns: final outcome after up to 3 retries.}
\label{tab:repair-transitions}
\small
\begin{tabular}{@{}lrrrrl@{}}
\toprule
Initial Failure & $\to$PASS & Persist. & Regr. & Total & Repair Rate \\
\midrule
BUILD\_FAIL & 72 & 274 & -- & 346 & 20.8\% [16.9\%, 25.4\%] \\ % src: paper_data.json > self_repair > per_initial_failure > BUILD_FAIL: 72/346=0.2081; CI [0.1686,0.2540]
RUN\_FAIL & 17 & 131 & -- & 148 & 11.5\% [7.3\%, 17.6\%] \\ % src: paper_data.json > self_repair > per_initial_failure > RUN_FAIL: 17/148=0.1149; CI [0.0730,0.1763]
VERIFY\_FAIL & 8 & 20 & -- & 28 & 28.6\% [15.2\%, 47.1\%] \\ % src: paper_data.json > self_repair > per_initial_failure > VERIFY_FAIL: 8/28=0.2857; CI [0.1525,0.4706]
EXTR\_FAIL & 15 & 12 & -- & 27 & 55.6\% [37.3\%, 72.4\%] \\ % src: paper_data.json > self_repair > per_initial_failure > EXTRACTION_FAIL: 15/27=0.5556; CI [0.3731,0.7241] (UNCHANGED)
UNKNOWN & 0 & 1 & -- & 1 & 0.0\% \\ % src: paper_data.json > self_repair > per_initial_failure > UNKNOWN: 0/1
\bottomrule
\end{tabular}
\end{table}


% src: paper_data.json > primary_campaign > self_repair (all values verified 2026-04-05)
% src: accounting: 160+112+392+7+39(partial_repair)=710; 160+112=272=Total PASS; 392+7+39=438=non-PASS
% src: 160/710=22.5%; 112/710=15.8%; 272/710=38.3%; (272-160)/160=70.0%; 7/710=1.0%
\begin{table}[htbp]
\centering
\caption{Self-repair effectiveness (all directions, all levels, Qwen~3.5; GPT-4.1~mini pending).}
\label{tab:self-repair}
\small
\begin{tabular}{@{}lrrl@{}}
\toprule
Metric & Qwen 3.5 & GPT-4.1 mini & Combined \\
\midrule
Total tasks & 710 & \tbd & \tbd \\
First-attempt PASS & 160 (22.5\%) & \tbd & \tbd \\ % src: paper_data.json > self_repair > first_attempt_pass=160, 160/710=0.2254
Repaired (attempt 2--4) & 112 (15.8\%) & \tbd & \tbd \\ % src: paper_data.json > self_repair > repaired=112, 112/710=0.1577
Total PASS & 272 (38.3\%) & \tbd & \tbd \\ % src: paper_data.json > primary_campaign > overall > pass=272, 272/710=0.3831
Relative improvement & +70.0\% & \tbd & \tbd \\ % src: (272-160)/160=112/160=0.70
Persistent fail & 392 & \tbd & \tbd \\ % src: paper_data.json > self_repair > persistent_fail=392
Regression & 7 (1.0\%) & \tbd & \tbd \\ % src: paper_data.json > self_repair > regression=7, 7/710=0.0099
\bottomrule
\end{tabular}
\end{table}


% src: paper_data.json > primary_campaign > augmentation > all_directions + cuda_to_omp_balanced (all 10 cells verified 2026-04-05)
\begin{table}[htbp]
\centering
\caption{LLM pass rates across augmentation levels. All directions ($n=142$ per level) and balanced CUDA-to-OpenMP subset ($n=24$ per level).}
\label{tab:augmentation-rates}
\small
\begin{tabular}{@{}cccc@{}}
\toprule
Level & \makecell{Qwen 3.5\\(all dirs, $n$=142)} & \makecell{Qwen 3.5\\(C$\to$OMP, $n$=24)} & \makecell{GPT-4.1\\mini} \\
\midrule
L0 & 40.1\% [32.4\%, 48.4\%] & 66.7\% [46.7\%, 82.0\%] & \tbd \\ % src: paper_data.json > augmentation > all_directions.L0: 57/142; cuda_to_omp_balanced.L0: 16/24
L1 & 37.3\% [29.8\%, 45.5\%] & 58.3\% [38.8\%, 75.5\%] & \tbd \\ % src: paper_data.json > augmentation > all_directions.L1: 53/142; cuda_to_omp_balanced.L1: 14/24
L2 & 40.1\% [32.4\%, 48.4\%] & 70.8\% [50.8\%, 85.1\%] & \tbd \\ % src: paper_data.json > augmentation > all_directions.L2: 57/142; cuda_to_omp_balanced.L2: 17/24
L3 & 39.4\% [31.8\%, 47.6\%] & 58.3\% [38.8\%, 75.5\%] & \tbd \\ % src: paper_data.json > augmentation > all_directions.L3: 56/142; cuda_to_omp_balanced.L3: 14/24
L4 & 34.5\% [27.2\%, 42.6\%] & 66.7\% [46.7\%, 82.0\%] & \tbd \\ % src: paper_data.json > augmentation > all_directions.L4: 49/142; cuda_to_omp_balanced.L4: 16/24
\bottomrule
\end{tabular}
\end{table}


\begin{figure}[htbp]
\centering
\includegraphics[width=\columnwidth]{f7_augmentation_robustness.pdf}
\caption{Augmentation robustness. X-axis: augmentation levels L0--L4. Y-axis: pass rate (\%). Lines per model. Rodinia CUDA-to-OpenMP, seed=42.}
\label{fig:augmentation}
\end{figure}


% src: paper_data.json > passk_campaign > passk_estimates (142 direction-kernel pairs, 3 samples each = 426 tasks)
% src: 103 hard_fail + 0 noisy + 17 always_pass + 22 partial = 142 verified
% src: 103/142=72.5%, 22/142=15.5%, 17/142=12.0% verified
% src: macro pass@1=19.7%, macro pass@3=27.5% verified; greedy pass rate=38.3% = primary campaign rate (paper_data.json > primary_campaign > overall > pass_rate=0.3831)
\begin{table*}[htbp]
\centering
\caption{pass@$k$ results (L0, temperature=0.7, 3 samples, Qwen~3.5; GPT-4.1~mini pending).}
\label{tab:pass-at-k}
\small
\begin{tabular}{@{}lrrrccc@{}}
\toprule
Model & Greedy rate (T=0) & pass@1 (T=0.7) & pass@3 (T=0.7) & Hard Fail & Noisy Fail & Always Pass \\
\midrule
Qwen 3.5 397B-A17B & 38.3\% (primary) & 19.7\% (macro) & 27.5\% (macro) & 103/142 (72.5\%) & 0/142 (0.0\%) & 17/142 (12.0\%) \\ % src: paper_data.json > primary_campaign > overall > pass_rate=0.3831 -> 38.3% greedy pass rate; passk unchanged
GPT-4.1 mini & \tbd & \tbd & \tbd & \tbd & \tbd & \tbd \\
\bottomrule
\end{tabular}
\end{table*}


\begin{figure}[htbp]
\centering
\includegraphics[width=\columnwidth]{f5_pass_at_k_by_direction.pdf}
\caption{pass@$k$ by translation direction.}
\label{fig:pass-at-k}
\end{figure}


\begin{figure}[htbp]
\centering
\includegraphics[width=\columnwidth]{f6_xsbench_comparison.pdf}
\caption{XSBench cross-granularity comparison. ParBench kernel-level vs.\ ParEval-Repo repository-level pass rates on the same computational kernel.}
\label{fig:xsbench}
\end{figure}


% src: paper_data.json > primary_campaign > augmentation > cochran_armitage: z=-0.0 (display 0.0), p_value=1.0
% src: paper_data.json > primary_campaign > direction_asymmetry: McNemar values cross-checked with S6.6 tab:direction-rates
% src: Bonferroni alpha = 0.05/4 = 0.0125 (4 McNemar tests incl. omp\_target)
\begin{table*}[htbp]
\centering
\caption{Statistical summary. All confidence intervals are Wilson 95\% CIs.}
\label{tab:stats-summary}
\small
\begin{tabular}{@{}lllll@{}}
\toprule
Test & Statistic & $p$-value & Effect Size & Interpretation \\
\midrule
Model comparison (chi-squared) & \tbd & \tbd & Cramer's V = \tbd & \tbd \\
Qwen augmentation (Cochran-Armitage) & $z = 0.0$ & $p = 1.0$ & -- & No monotonic trend; null supports robustness \\ % src: paper_data.json > augmentation > cochran_armitage: z=-0.0 (display 0.0), p_value=1.0, n_kernels=24
GPT-4.1~mini augmentation (Cochran-Armitage) & \tbd & \tbd & -- & \tbd \\
Direction: cuda-omp (McNemar) & $n=24$, exact & $p = 0.6875$ & Cohen's $h = -0.17$ & No significant asymmetry ($\alpha=0.0125$) \\ % src: paper_data.json > direction_asymmetry > omp-to-cuda vs cuda-to-omp: n=24, p=0.6875, h=-0.1724
Direction: cuda-opencl (McNemar) & $n=20$, exact & $p = 0.6875$ & Cohen's $h = 0.28$ & No significant asymmetry ($\alpha=0.0125$) \\ % src: paper_data.json > direction_asymmetry > cuda-to-opencl vs opencl-to-cuda: n=20, p=0.6875, h=0.2838
Direction: omp-opencl (McNemar) & $n=18$, exact & $p = 1.0$ & Cohen's $h = 0.12$ & No significant asymmetry ($\alpha=0.0125$) \\ % src: paper_data.json > direction_asymmetry > opencl-to-omp vs omp-to-opencl: n=18, p=1.0, h=0.1157
\bottomrule
\end{tabular}
\end{table*}

% --- Full Per-Kernel Table (complete 31-row version of condensed main-text table) ---


% src: paper_data.json > primary_campaign > by_kernel (all 31 rows verified 2026-04-05)
% src: row sums all match Total; total PASS across kernels = 272; total tasks = 710
% src: 31 evaluated kernels = 35 corpus kernels - 4 with all specs KNOWN_FAIL/phantom (kmeans, mummergpu, hybridsort, huffman)
% src: spot-checked vs raw JSONs: hotspot(21P/4R/5V=30), stencil1d(15P=15), xsbench(20B/10R=30)
\begin{table*}[htbp]
\centering
\caption{Full per-kernel greedy pass rates across all 31 evaluation kernels (Qwen~3.5 397B-A17B, 710 primary campaign tasks, 31~kernels across 5 suites). KNOWN\_FAIL source specs excluded.}
\label{tab:per-kernel-full}
\footnotesize
\begin{tabular}{@{}llrrrrl@{}}
\toprule
Suite & Kernel & Total & PASS & Fail & Rate & 95\% Wilson CI \\
\midrule
% src: paper_data.json > primary_campaign > by_kernel > stencil1d: pass=15, total=15, rate=1.0
HeCBench & stencil1d & 15 & 15 & 0 & 100.0\% & [79.6\%, 100.0\%] \\
% src: paper_data.json > primary_campaign > by_kernel > floydwarshall: pass=16, total=20, rate=0.8
HeCBench & floydwarshall & 20 & 16 & 4 & 80.0\% & [58.4\%, 91.9\%] \\
% src: paper_data.json > primary_campaign > by_kernel > heat2d: pass=15, total=20, rate=0.75
HeCBench & heat2d & 20 & 15 & 5 & 75.0\% & [53.1\%, 88.8\%] \\
% src: paper_data.json > primary_campaign > by_kernel > iso2dfd: pass=15, total=20, rate=0.75
HeCBench & iso2dfd & 20 & 15 & 5 & 75.0\% & [53.1\%, 88.8\%] \\
% src: paper_data.json > primary_campaign > by_kernel > hotspot: pass=21, total=30, rate=0.7
Rodinia & hotspot & 30 & 21 & 9 & 70.0\% & [52.1\%, 83.3\%] \\
% src: paper_data.json > primary_campaign > by_kernel > nn: pass=7, total=10, rate=0.7
Rodinia & nn & 10 & 7 & 3 & 70.0\% & [39.7\%, 89.2\%] \\
% src: paper_data.json > primary_campaign > by_kernel > hotspot3d: pass=19, total=30, rate=0.6333
Rodinia & hotspot3d & 30 & 19 & 11 & 63.3\% & [45.5\%, 78.1\%] \\
% src: paper_data.json > primary_campaign > by_kernel > jacobi: pass=6, total=10, rate=0.6
HeCBench & jacobi & 10 & 6 & 4 & 60.0\% & [31.3\%, 83.2\%] \\
% src: paper_data.json > primary_campaign > by_kernel > md: pass=6, total=10, rate=0.6
HeCBench & md & 10 & 6 & 4 & 60.0\% & [31.3\%, 83.2\%] \\
% src: paper_data.json > primary_campaign > by_kernel > nqueen: pass=6, total=10, rate=0.6
HeCBench & nqueen & 10 & 6 & 4 & 60.0\% & [31.3\%, 83.2\%] \\
% src: paper_data.json > primary_campaign > by_kernel > page-rank: pass=6, total=10, rate=0.6
HeCBench & page-rank & 10 & 6 & 4 & 60.0\% & [31.3\%, 83.2\%] \\
% src: paper_data.json > primary_campaign > by_kernel > particlefilter: pass=17, total=30, rate=0.5667
Rodinia & particlefilter & 30 & 17 & 13 & 56.7\% & [39.2\%, 72.6\%] \\
% src: paper_data.json > primary_campaign > by_kernel > bfs: pass=16, total=30, rate=0.5333
Rodinia & bfs & 30 & 16 & 14 & 53.3\% & [36.1\%, 69.8\%] \\
% src: paper_data.json > primary_campaign > by_kernel > cfd: pass=16, total=30, rate=0.5333
Rodinia & cfd & 30 & 16 & 14 & 53.3\% & [36.1\%, 69.8\%] \\
% src: paper_data.json > primary_campaign > by_kernel > lud: pass=15, total=30, rate=0.5; Wilson CI [33.1539%, 66.8461%]
Rodinia & lud & 30 & 15 & 15 & 50.0\% & [33.2\%, 66.8\%] \\
% src: paper_data.json > primary_campaign > by_kernel > nw: pass=15, total=30, rate=0.5; Wilson CI [33.1539%, 66.8461%]
Rodinia & nw & 30 & 15 & 15 & 50.0\% & [33.2\%, 66.8\%] \\
% src: paper_data.json > primary_campaign > by_kernel > pathfinder: pass=15, total=30, rate=0.5; Wilson CI [33.1539%, 66.8461%]
Rodinia & pathfinder & 30 & 15 & 15 & 50.0\% & [33.2\%, 66.8\%] \\
\midrule
% src: paper_data.json > primary_campaign > by_kernel > backprop: pass=13, total=30, rate=0.4333
Rodinia & backprop & 30 & 13 & 17 & 43.3\% & [27.4\%, 60.8\%] \\
% src: paper_data.json > primary_campaign > by_kernel > srad: pass=12, total=30, rate=0.4
Rodinia & srad & 30 & 12 & 18 & 40.0\% & [24.6\%, 57.7\%] \\
% src: paper_data.json > primary_campaign > by_kernel > scan: pass=5, total=15, rate=0.3333
HeCBench & scan & 15 & 5 & 10 & 33.3\% & [15.2\%, 58.3\%] \\
% src: paper_data.json > primary_campaign > by_kernel > mixbench: pass=8, total=30, rate=0.2667; Wilson CI [14.1825%, 44.4483%]
mixbench & mixbench & 30 & 8 & 22 & 26.7\% & [14.2\%, 44.4\%] \\
% src: paper_data.json > primary_campaign > by_kernel > bptree: pass=4, total=30, rate=0.1333
Rodinia & bptree & 30 & 4 & 26 & 13.3\% & [5.3\%, 29.7\%] \\
% src: paper_data.json > primary_campaign > by_kernel > lavamd: pass=2, total=30, rate=0.0667
Rodinia & lavamd & 30 & 2 & 28 & 6.7\% & [1.8\%, 21.3\%] \\
% src: paper_data.json > primary_campaign > by_kernel > streamcluster: pass=2, total=30, rate=0.0667
Rodinia & streamcluster & 30 & 2 & 28 & 6.7\% & [1.8\%, 21.3\%] \\
\midrule
% src: paper_data.json > primary_campaign > by_kernel > convolution1d: pass=0, total=10, rate=0.0
HeCBench & convolution1d & 10 & 0 & 10 & 0.0\% & [0.0\%, 27.8\%] \\
% src: paper_data.json > primary_campaign > by_kernel > dwt2d: pass=0, total=10, rate=0.0
Rodinia & dwt2d & 10 & 0 & 10 & 0.0\% & [0.0\%, 27.8\%] \\
% src: paper_data.json > primary_campaign > by_kernel > gaussian: pass=0, total=10, rate=0.0
Rodinia & gaussian & 10 & 0 & 10 & 0.0\% & [0.0\%, 27.8\%] \\
% src: paper_data.json > primary_campaign > by_kernel > heartwall: pass=0, total=30, rate=0.0; Wilson CI [0.0%, 11.3517%]
Rodinia & heartwall & 30 & 0 & 30 & 0.0\% & [0.0\%, 11.4\%] \\
% src: paper_data.json > primary_campaign > by_kernel > myocyte: pass=0, total=30, rate=0.0; Wilson CI [0.0%, 11.3517%]
Rodinia & myocyte & 30 & 0 & 30 & 0.0\% & [0.0\%, 11.4\%] \\
% src: paper_data.json > primary_campaign > by_kernel > rsbench: pass=0, total=30, rate=0.0; Wilson CI [0.0%, 11.3517%]
RSBench & rsbench & 30 & 0 & 30 & 0.0\% & [0.0\%, 11.4\%] \\
% src: paper_data.json > primary_campaign > by_kernel > xsbench: pass=0, total=30, rate=0.0; Wilson CI [0.0%, 11.3517%]
XSBench & xsbench & 30 & 0 & 30 & 0.0\% & [0.0\%, 11.4\%] \\
\bottomrule
\end{tabular}
\end{table*}

